\section{From local jets to a global coefficient remainder}\label{cms:sec:global}
The next argument is the analytic step that distinguishes a coefficient theorem from a radial heuristic. We fix a target order and verify the boundary regularity required for Fourier integration by parts.

\subsection{Choice of finite cutoffs}
Fix $K\geq2$, and set
\begin{equation}\label{cms:eq:cutoffs}
M=K+2,\qquad L=K+6,\qquad d=L-1=K+5=M+3.
\end{equation}
By Lemma~\ref{cms:lem:tail}, $Q_L$ has bounded continuous complex derivatives through order $d$. Let $Z_L$ be the finite set of roots of unity with exact order at most $L$. This is the singular set of the finite factor $P_L$ on the unit circle. It is not the full singular set of $F$.

At each nontrivial $\zeta\in Z_L$, use \eqref{cms:eq:rootform}. Expand its analytic factor in Taylor series, replace $Q_L$ by its Taylor polynomial of degree $d-1$, and expand the exponential of $-B_{q,L}\log s$. Retain every nonanalytic term that becomes
\begin{equation}\label{cms:eq:jetmonomial}
 c_{\zeta,p,h}\left(1-\frac z\zeta\right)^p
 \log^h\left(1-\frac z\zeta\right),\qquad
 0\leq p\leq M-1,\quad h\geq1,
\end{equation}
after the analytic change from $s$ to $t=1-z/\zeta$. There are only finitely many such terms. The minimum degree in $B_{q,L}$ is $q-1\geq1$, so a power $(\log s)^h$ from the exponential has degree at least $h(q-1)$ before the possible extra zero. This ensures finiteness without assigning a cutoff in logarithm degree separately.

The coefficients $c_{\zeta,p,h}$ are defined by this finite-jet operation. They do not depend on increasing $L$ once sufficient differentiability and Taylor depth are available: two valid expansions with controlled higher-order remainders have identical lower-order nonanalytic jets. Their values can include convergent derivatives of $Q_L$, rather than only elementary constants. The Gamma products are their leading amplitudes.

\subsection{Remainders at the nontrivial selected roots}
We record why each step in this local construction remains valid after $M$ derivatives. For a fixed integer $p\geq M$ and $h\geq0$,
\begin{equation}\label{cms:eq:monomialderivatives}
\frac{\mathrm d^M}{\mathrm ds^M}\bigl(s^p\log^h s\bigr)
=O\!\left(|s|^{p-M}(1+|\log s|)^h\right)
\end{equation}
on the relevant punctured neighborhood. Derivatives of lower orders have the analogous extra positive powers. In particular, when $p=M$, the $M$th derivative is locally integrable on a boundary arc, and derivatives through order $M-1$ tend to zero for this remainder.

The exponential in \eqref{cms:eq:rootform} is expanded to enough terms that its first omitted power has $s$-degree at least $M$. Its usual Taylor formula, applied to the small quantity $B_{q,L}(s)\log s$, yields the same differentiated bounds. Truncating analytic factors is equally harmless after choosing their Taylor depth beyond $M$.

The only nonanalytic object not represented by a convergent power-log expansion is $Q_L$. Lemma~\ref{cms:lem:taylor} gives, in the $s$ coordinate,
\begin{equation}\label{cms:eq:Qremainder}
\left(Q_L(\zeta\ee^{-s})-T(s)\right)^{(h)}
=O(|s|^{d-h})\qquad(0\leq h\leq d),
\end{equation}
for a Taylor polynomial $T$ of degree $d-1$. Derivatives of $P_L$ of order $v\leq M$ consist of bounded analytic contributions and terms bounded by
\[
O\!\left(|s|^{\varepsilon_\zeta+q-1-v}
 (1+|\log s|)^{c_v}\right)
\]
for fixed finite exponents $c_v$. The Leibniz rule applied to $P_L(Q_L-T)$ now gives positive powers of $|s|$ after $M$ derivatives. For the singular part the exponent is at least
\[
d+\varepsilon_\zeta+q-1-M\geq4,
\]
and the bounded analytic part gives at least $d-M=3$. All these terms are locally integrable, with lower derivatives continuous at the root. This explicitly checks the tail remainder rather than treating it as an analytic function across the boundary.

The analytic change $s=-\log(1-t)$ preserves these bounds: $s/t$ is analytic and nonzero near zero, and
\[
\log s=\log t+\log(s/t)
\]
adds only an analytic factor. Differentiation along the boundary angular coordinate also preserves the local integrability and continuity properties, since it is a nonsingular smooth coordinate change.

\subsection{The remainder at the positive root}
At $1$, equation \eqref{cms:eq:positiveform} gives
\[
F(\ee^{-s})=s^{-1}U_L(s)Q_L(\ee^{-s})
 \exp\bigl(D_L(s)\log s\bigr),\qquad D_L(s)=O(s^L).
\]
The first logarithmic term, including the pole, has degree $L-1$. Its $M$th derivative is bounded by
\begin{equation}\label{cms:eq:positiveerror}
O\!\left(|s|^{L-1-M}(1+|\log s|)^c\right)
=O\!\left(|s|^3(1+|\log s|)^c\right).
\end{equation}
The remainder from replacing $Q_L$ by its degree-$(d-1)$ Taylor polynomial and dividing by $s$ has $M$th derivative
\begin{equation}\label{cms:eq:positiveQerror}
O(|s|^{d-M-1})=O(|s|^2).
\end{equation}
After the pole $A/(1-z)$ is subtracted, everything else up to these remainders is analytic near $1$. Equations \eqref{cms:eq:positiveerror}--\eqref{cms:eq:positiveQerror} therefore give the same required boundary regularity there.

\subsection{A single analytic subtraction in the disk}
For each $\zeta\in\partial\D$, the function $\log(1-z/\zeta)$ has its usual analytic branch in $\D$, since $1-z/\zeta$ has positive real part there. Let $J_\zeta$ be the finite sum \eqref{cms:eq:jetmonomial} at the nontrivial root $\zeta$. Define
\begin{equation}\label{cms:eq:G}
G(z)=F(z)-\frac{A}{1-z}-\sum_{\zeta\in Z_L\setminus\{1\}}J_\zeta(z).
\end{equation}
Every subtraction is analytic in the open disk. At a different selected root, a fixed $J_\zeta$ is locally analytic. Subtracting it may change ordinary Taylor coefficients but cannot change a nonanalytic jet at that other root. Thus all of the local estimates above hold simultaneously for $G$.

Away from the finite set $Z_L$, the finite factor $P_L$ continues analytically and $Q_L$ has $d$ bounded boundary derivatives. The high-order roots not in $Z_L$ are included in that smooth tail. They do not create an unexamined exceptional set in the argument. A finite covering of the circle by selected-root neighborhoods and their complement gives a global conclusion.

\begin{lemma}[Boundary regularity and Fourier estimate]\label{cms:lem:Fourier}
The function $G$ in \eqref{cms:eq:G} extends continuously to $\overline\D$. Its boundary trace $g(\theta)=G(\ee^{\ii\theta})$ has continuous derivatives through order $M-1$, these derivatives are absolutely continuous, and $g^{(M)}\in L^1[0,2\pi]$. Consequently
\begin{equation}\label{cms:eq:globalremainder}
[z^m]G(z)=O_K(m^{-M})=O_K(m^{-K-2}).
\end{equation}
\end{lemma}
\begin{proof}
The local remainders have continuous derivatives through order $M-1$ by the estimates just established. On each punctured arc they possess an ordinary $M$th derivative that is integrable, since all powers of $|\log |\theta-\theta_0||$ are locally integrable. The lower-order derivatives have matching limits from either side. The fundamental theorem of calculus on truncated arcs, followed by the integrable limit, proves absolute continuity across the removed point. A finite partition then proves global absolute continuity. This also verifies the asserted regularity of all lower derivatives.

Since $G$ is continuous on the compact closed disk, its radial boundary values converge uniformly. Cauchy's coefficient formula therefore identifies its Taylor coefficients with the Fourier coefficients of $g$. Repeated integration by parts, with periodic boundary terms canceling, gives for $m\geq1$
\begin{equation}\label{cms:eq:Fourier}
[z^m]G(z)=\frac{(\ii m)^{-M}}{2\pi}
\int_0^{2\pi}g^{(M)}(\theta)\ee^{-\ii m\theta}\dd\theta.
\end{equation}
The $L^1$ norm proves \eqref{cms:eq:globalremainder}. The Riemann--Lebesgue lemma would strengthen this to a little-oh estimate at the same power, but is not needed.
\end{proof}

\subsection{Relationship to the published hybrid theorem}
This is a direct boundary proof within the classical Darboux and hybrid framework \cite{FFGPP}. To invoke their Theorem~2 literally instead, one can choose the larger cutoff $L=2K+6$ and boundary regularity $s=2K+5$. The singular factor has global order $a=-1$, because \eqref{cms:eq:globalbound} holds and $Q_L^{-1}$ is bounded. Their exponent is
\[
u_0=\left\lfloor\frac{s+\lfloor a\rfloor}{2}\right\rfloor=K+2.
\]
The finite polylogarithms continue into a common indented domain at the finitely many roots and have arbitrarily deep power-log expansions, so the remaining hypotheses can be met by choosing their local expansion depth $t_0>u_0$. Their conclusion gives the required coefficient remainder. The smaller cutoff $L=K+6$ in \eqref{cms:eq:cutoffs} is justified by the direct derivative argument above; it is not obtained by silently substituting into the published exponent formula.
